\subsection{Conceptos de trabajo del sistema}
\label{subsec:conceptos_trabajo}

En este apartado se detallan las reglas operativas con las que trabaja la página web para evitar la pérdida de archivos y asegurar que los acuerdos se cumplan puntualmente.

\subsubsection{Control de versiones de documentos}
El control de versiones consiste en guardar y ordenar copias numeradas de un mismo archivo (Versión 1, Versión 2, etc.) cada vez que se le agregan avances o correcciones \parencite{chacon2014progit}. En la página web desarrollada se aplica una regla obligatoria de preservación: ninguna entrega puede sobreescribir ni borrar los archivos pasados. Gracias a esto, tanto el alumno como el asesor pueden consultar en cualquier momento los borradores anteriores para revisar comentarios previos sin perder información.

\subsubsection{Uso exclusivo de archivos en formato PDF}
El formato PDF (\textit{Portable Document Format}) es un estándar internacional formalizado mediante la norma técnica ISO 32000-2 \parencite{iso32000}. La página web solicita que los borradores se suban exclusivamente en este formato porque asegura que los textos, tipos de letra, figuras, fórmulas matemáticas y tablas se vean exactamente igual en cualquier computadora, teléfono o sistema operativo, evitando que el documento se mueva de lugar como sucede frecuentemente en procesadores editables.

\subsubsection{Minutas digitales de asesoría}
Una minuta es un registro escrito que se llena al terminar una junta o reunión para dejar asentados los temas que se platicaron, los acuerdos tomados y las fechas pactadas de entrega \parencite{pmi2021}. La página web incluye un formulario específico donde el asesor escribe estos acuerdos al término de cada sesión, de modo que ambas partes tengan muy claro qué tareas deben entregar antes de la siguiente cita de trabajo.

\subsubsection{Regla de aviso por inactividad mayor a 30 días}
Para evitar que los estudiantes se queden rezagados en el posgrado, la página web integra un reloj automático que compara la fecha de hoy contra el último día en que el tesista subió un avance en PDF o registró una minuta de asesoría. Si la cuenta llega a más de 30 días naturales sin ninguna actividad, el sistema coloca una alerta visual en el expediente del alumno para que el profesor y la coordinación se comuniquen con él a tiempo y le ayuden a retomar su tesis.